\medskip

\emph{Cet exercice est un questionnaire à choix multiples. \\
Pour chacune des six questions suivantes, une seule des quatre réponses proposées est exacte.\\
Une réponse fausse, une réponse multiple ou l'absence de réponse à une question ne rapporte ni n'enlève de point.\\
Pour répondre, indiquer sur la copie le numéro de la question et la lettre de la réponse choisie.\\
Aucune justification n'est demandée.}

\medskip

\begin{enumerate}
\item On considère la fonction $f$ définie et dérivable sur $]0~;~+\infty[$ par:
\[f(x) = x \ln(x) - x + 1.\]

Parmi les quatre expressions suivantes, laquelle est celle de la fonction dérivée de $f$?

\begin{center}
\begin{tabularx}{\linewidth}{|*{4}{X|}}\hline
\textbf{a.~~} $\ln (x)$&\textbf{b.~~}$\dfrac{1}{x} - 1$&\textbf{c.~~} $\ln (x) - 2$&\textbf{d.~~}$\ln (x) - 1$\rule[-3mm]{0mm}{9mm}\\ \hline
\end{tabularx}
\end{center}

\item On considère la fonction $g$ définie sur $]0~;~+\infty[$ par $g(x) = x^2[1 - \ln (x)]$.

 Parmi les quatre affirmations suivantes, laquelle est correcte ?
 
\begin{center}
\begin{tabularx}{\linewidth}{|*{4}{X|}}\hline
\textbf{a.~~}$\displaystyle\lim_{x \to 0} g(x) = +\infty$&\textbf{b.~~} 
$\displaystyle\lim_{x \to 0} g(x) = - \infty$&\textbf{c.~~} $\displaystyle\lim_{x \to 0} g(x) = 0$&\textbf{d.~~} La fonction $g$ n'admet pas de limite en 0.\\ \hline
\end{tabularx}
\end{center}

\item On considère la fonction $f$ définie sur $\R$ par $f(x) = x^3 - 0,9x^2 -0,1x$. Le nombre de solutions de l'équation $f(x) = 0$ sur $\R$ est :

\begin{center}
\begin{tabularx}{\linewidth}{|*{4}{X|}}\hline
\textbf{a.~~} $0$&\textbf{b.~~}$1$&\textbf{c.~~} $2$&\textbf{d.~~}$3$\\ \hline
\end{tabularx}
\end{center}

\item Si $H$ est une primitive d'une fonction $h$ définie et continue sur $\R$, 
et si $k$ est la fonction définie sur $\R$ par $k(x) = h(2x)$, 
alors, une primitive $K$ de $k$ est définie sur $\R$ par :

\begin{center}
\begin{tabularx}{\linewidth}{|*{4}{X|}}\hline
\textbf{a.~~} $K(x) =H(2x)$&\textbf{b.~~}$K(x) =2H(2x)$&\textbf{c.~~} $ K(x) =\dfrac{1}{2}H(2x)$&\textbf{d.~~}$K(x) =2H(x)$\rule[-3mm]{0mm}{9mm}\\ \hline
\end{tabularx}
\end{center}

\item L'équation réduite de la tangente au point d'abscisse 1 de la courbe de la fonction $f$ définie sur $\R$ par $f(x) = x\text{e}^x$ est:

\begin{center}
\begin{tabularx}{\linewidth}{|*{4}{X|}}\hline
\textbf{a.~~} $y = \text{e}x + \text{e}$&\textbf{b.~~}$y =2\text{e}x - \text{e}$&\textbf{c.~~} $y = 2\text{e}x +  \text{e}$&\textbf{d.~~}$y = \text{e}x$\\ \hline
\end{tabularx}
\end{center}

\item Les nombres entiers $n$ solutions de l'inéquation $(0,2)^n < 0,001$ sont tous les
nombres entiers $n$ tels que :

\begin{center}
\begin{tabularx}{\linewidth}{|*{4}{X|}}\hline
\textbf{a.~~} $n \leqslant 4$&\textbf{b.~~}$n\leqslant 5$&\textbf{c.~~} $n \geqslant 4$ &\textbf{d.~~}$n \geqslant 5$\\ \hline
\end{tabularx}
\end{center}

\end{enumerate}

\bigskip
